\chapter{Experimental Evaluation and Results}\label{ch:results}

\section{Experimental Evaluation Overview}
The evaluation of TrustLayer-AI was conducted across distinct empirical phases spanning model diagnostics, information retrieval ablation, natural language grounding validation, database storage parity checks, data pipeline integrity verification, and master backend provenance testing. 

Rather than assuming ideal system execution from inception, the project followed an iterative, evidence-driven evaluation paradigm. Empirical findings from early stages directly triggered major architectural remediations, including overhauling interaction modeling, switching to Reciprocal Rank Fusion (RRF), abandoning high-latency SHAP explainers in favor of analytical feature matching, and migrating file-based stores to an enterprise PostgreSQL 17.6 + \texttt{pgvector} database.

\section{Dataset and Data Processing Results}
Data acquisition and cleaning yielded verified statistical metrics across Delhi NCR properties:
\begin{itemize}
    \item \textbf{Raw Property Acquisition}: 1,661 distinct hotel entities collected via the Google Places Details API across 15 major area clusters.
    \item \textbf{Raw Text Reviews}: 8,305 raw text review records acquired, reflecting the Google Places Detail API request limit of 5 reviews per property.
    \item \textbf{Cleaned Metadata and Reviews}: 1,661 hotels preserved after geographic bounding-box filtering ($28.40^\circ\text{N} \le \text{Lat} \le 28.88^\circ\text{N}$, $76.84^\circ\text{E} \le \text{Lng} \le 77.35^\circ\text{E}$). 1,618 hotels contained non-empty text reviews, while 43 hotels had 0 reviews.
    \item \textbf{Canonical Handoff Dataset}: `\texttt{final\_hotel\_dataset.csv}` established with 1,661 hotels and 26 engineered feature columns (SHA-256: `\texttt{eca959c788d9535feea5ed1b11efb249988c60b6bfb7b4bb41ef33e50bad1e2e}`).
    \item \textbf{Synthetic User Profiles}: 500 demographic user preference profiles generated with budget, area, travel purpose, and aspect priority attributes.
    \item \textbf{Review Evidence Segmentation}: 7,910 text chunks generated across Chunk Types A (Profile), B (Aspects), C (Positive Evidence), D (Negative Evidence), and E (Recommender Signals).
\end{itemize}

Geographic clustering confirmed that hotels are heavily concentrated in central New Delhi (Paharganj and Karol Bagh, 216 properties), the Indira Gandhi International Airport corridor (Mahipalpur, 56 properties), and Gurugram (39 properties). The 100\% missing \texttt{price\_level} attribute returned by the Places API was successfully mitigated by engineering a \texttt{budget\_category} proxy feature based on area cluster tariffs and rating tiers.

\section{Sentiment Analysis and ABSA Results}
NLP evaluation validated sentence-level sentiment extraction and multi-dimensional aspect scoring:
\begin{itemize}
    \item \textbf{DistilBERT Sentiment Polarity}: Sentence-level positive sentiment probabilities ($P_{\text{pos}}$) extracted via DistilBERT (`\texttt{distilbert-base-uncased-finetuned-sst-2-english}`) demonstrated a high Pearson correlation coefficient of **$r \approx 0.84$** against user star ratings, validating NLP sentiment extraction quality.
    \item \textbf{Aspect-Based Sentiment Distributions}: ABSA scores ($S_{\text{aspect}} \in [0, 100]$) evaluated across Cleanliness, Service, Location, Value for Money, and Staff Behavior revealed that **Cleanliness** exhibited the highest score variance across Delhi NCR properties, establishing it as the key aspect differentiator for hotel selection.
\end{itemize}

\section{Trust and Feature Engineering Results}
Feature engineering produced two orthogonal quality indicators:
\begin{enumerate}
    \item \textbf{Trust Score ($T_{\text{score}}$)}: Combining normalized star rating, DistilBERT sentiment probability, and review volume support produced a Gaussian distribution centered at $68.0$ (range $0$ to $100$).
    \item \textbf{Popularity Score ($P_{\text{score}}$)}: Log-transformed review volume yielded a steep power-law distribution.
\end{enumerate}

Evaluating the correlation between Trust Score and Popularity Score yielded **$r \approx 0.05$**, proving that Trust and Popularity represent independent evaluation signals. A hotel with modest review volume can achieve a high Trust Score if its guest reviews exhibit consistently high sentiment and aspect ratings.

\section{Initial Recommendation System Diagnostic Failure (Stage A)}
Offline evaluation of the initial recommender algorithms in Stage A (`\texttt{recommender\_diagnostics.md}`) uncovered severe performance degradation:
\begin{itemize}
    \item \textbf{Precision@10}: $0.002$
    \item \textbf{Recall@10}: $0.010$
    \item \textbf{NDCG@10}: $0.006$
\end{itemize}

These values were barely above random chance, triggering a formal **NO-GO** decision for downstream deployment. Diagnostic investigation identified three root causes:
\begin{enumerate}
    \item \textbf{Extreme Matrix Sparsity}: The initial user-item interaction matrix (500 users, 1,379 hotels, 5,000 interactions) exhibited a **99.27\% matrix sparsity** (~3.6 interactions/hotel). Collaborative filtering via SVD underfit heavily, failing to learn latent factors and reverting to predicting global item biases.
    \item \textbf{Score Calibration Mismatch}: Content-based cosine similarities clustered densely between $0.8$ and $0.9$, while SVD predicted ratings spread from $1.0$ to $5.0$. Grid-search optimization over the linear hybrid model ($\alpha \cdot S_{\text{CF}} + (1-\alpha) \cdot S_{\text{CB}}$) selected $\alpha = 1.0$, completely disabling the content-based component.
    \item \textbf{Evaluation Script Reporting Flaw}: Auditing `\texttt{evaluate\_recommenders.py}` revealed that narrative text claiming Content-Based filtering outperformed Collaborative Filtering on cold-start users was statically hardcoded. User-level 70/30 chronological splitting over fixed 10-interaction profiles left exactly **0 cold-start users** in the test set.
\end{enumerate}

\section{Recommendation System Remediation and RRF Results (Stage A.1)}
To resolve these failures, two major remediations were executed in Stage A.1 (`\texttt{recommender\_remediation_report.md}`):
\begin{enumerate}
    \item \textbf{V2 Synthetic Interaction Generation}: Regenerated user interactions (`\texttt{generate\_interactions\_v2.py}`) with realistic preference overlap (66\% budget match, 51\% area match) and power-law interaction frequencies.
    \item \textbf{Reciprocal Rank Fusion (RRF)}: Replaced linear score blending with Reciprocal Rank Fusion ($k=60$). RRF combines ordinal item ranks rather than raw prediction scores, eliminating score calibration mismatches.
\end{enumerate}

Re-evaluation demonstrated dramatic ranking quality recovery, elevating **$\text{NDCG}@10$ from $0.006$ to $> 0.12$**, balancing catalog coverage, and securing a formal **GO** decision.

\section{Explainability Evaluation}
Evaluating explainability mechanisms highlighted significant performance trade-offs:
\begin{itemize}
    \item \textbf{SHAP Evaluation}: Initial SHAP value calculations incurred severe computational latency ($> 1.5$ seconds per query) and produced dense, non-deterministic feature importance scores over RRF rank-fused outputs that failed human readability audits.
    \item \textbf{Analytical Feature-Matching Explainer}: Replacing SHAP with an analytical explainer (`\texttt{explainer.py}`) enabled direct computation of aspect alignment scores across Cleanliness, Service, Location, Value, and Staff Behavior in **$< 5$ ms**. Auditing over 100 sample user profiles passed 100\% of logic and readability checks.
\end{itemize}

\section{RAG Retrieval Evaluation and Ablation}
Hybrid retrieval evaluation over 150 benchmark queries (`\texttt{retrieval\_evaluation.md}`) validated dense vector retrieval, metadata filtering, and recommender reranking against success gates.

Table~\ref{tab:retrieval_evaluation} presents the retrieval performance across configuration stages.

\begin{table}[!ht]
\centering
\caption{RAG Retrieval Ablation Study across 150 Benchmark Queries}
\label{tab:retrieval_evaluation}
\begin{tabular}{lcccc}
\hline
\textbf{Retrieval Configuration} & \textbf{Precision@5} & \textbf{Recall@5} & \textbf{MRR} & \textbf{NDCG@5} \\
\hline
Semantic Vector Only & 0.68 & 0.65 & 0.70 & 0.71 \\
+ Metadata Filtering & 0.75 & 0.73 & 0.78 & 0.79 \\
+ Recommendation Reranking & 0.73 & 0.70 & 0.75 & 0.76 \\
\textbf{Full Hybrid Retrieval} & \textbf{0.81} & \textbf{0.79} & \textbf{0.83} & \textbf{0.85} \\
\hline
\textbf{Target Threshold} & $\ge 0.70$ & $\ge 0.70$ & $\ge 0.75$ & $\ge 0.75$ \\
\textbf{Status} & \textbf{PASS} & \textbf{PASS} & \textbf{PASS} & \textbf{PASS} \\
\hline
\end{tabular}
\end{table}

Ablation analysis proved that adding hard metadata filtering contributed $+0.07$ to Precision@5, while recommender reranking contributed an additional $+0.05$, achieving a final Precision@5 of **0.81**. The total retrieval latency breakdown was measured as: Embedding Generation (12.4 ms), ChromaDB Search (145.2 ms), Metadata Filtering (2.1 ms), Reranking (4.3 ms), and Context Assembly (1.8 ms), totaling **165.8 ms**.

\section{Grounding and Hallucination Evaluation}
Evaluating the RAG answer generation pipeline and `\texttt{GroundingValidator}` across 150 benchmark queries (`\texttt{stage\_d\_evaluation.md}`) yielded verified reliability metrics:
\begin{itemize}
    \item \textbf{Grounded Response Rate}: **96.7%** (Target $\ge 95\%$) — PASS
    \item \textbf{Hallucination Rate}: **1.3%** (Target $\le 2\%$) — PASS
    \item \textbf{Explanation Consistency}: **96.0%** (Target $\ge 95\%$) — PASS
    \item \textbf{Provenance Coverage}: **98.7%** (Target $\ge 95\%$) — PASS
    \item \textbf{Active Interceptions}: `\texttt{GroundingValidator}` actively intercepted **3 responses** containing fabricated amenity assertions (e.g., claiming a budget hotel featured a luxury spa), stripping ungrounded text and triggering context re-prompting.
    \item \textbf{Cache Performance}: Cache Miss Latency = $210.0$ ms | Cache Hit Latency = $12.0$ ms.
\end{itemize}

\section{ChromaDB $\rightarrow$ PostgreSQL/pgvector Migration Results (Stage 24.2)}
Stage 24.2 backfilled all 7,910 review chunks into PostgreSQL 17.6 `\texttt{pgvector}` storage (`\texttt{STAGE24.2\_PGVECTOR\_BACKFILL\_REPORT.md}`).

Controlled dual-backend verification demonstrated:
\begin{itemize}
    \item \textbf{Embedding Cosine Similarity Parity}: Sampling 100 random vectors across ChromaDB and `\texttt{pgvector}` yielded an **average cosine similarity of 1.0000** ($0.0000$ max absolute difference across 384 dimensions).
    \item \textbf{20-Query RAG Parity}: Evaluating 20 recommendation queries comparing legacy file stores against PostgreSQL achieved **20 / 20 (100.0%) Top-1 hotel match** and **100% Top-5 candidate overlap**.
\end{itemize}

\section{PostgreSQL Relational Integrity Results (Stage 24.5)}
Master provenance auditing (`\texttt{STAGE24.5\_COMPLETE\_BACKEND\_FORENSIC\_VERIFICATION\_REPORT.md}`) verified direct SQL query counts in the production PostgreSQL database (`\texttt{trustlayer\_db}`):
\begin{itemize}
    \item `\texttt{hotels}`: **1,661 rows**
    \item `\texttt{hotel\_locations}`: **1,661 rows**
    \item `\texttt{hotel\_scores}`: **1,661 rows**
    \item `\texttt{hotel\_sources}`: **1,661 rows**
    \item `\texttt{hotel\_amenities}`: **1,661 rows**
    \item `\texttt{embedding\_documents}`: **7,910 rows**
    \item `\texttt{domain\_events}`: **1,661 rows** (Transactional outbox events)
    \item `\texttt{ingestion\_records}`: **1,661 rows**
    \item `\texttt{ingestion\_runs}`: **1 row**
\end{itemize}

Relational checks confirmed **0 orphan records** across child tables and **0 duplicate primary keys**, validating complete legacy file detachment.

\section{Stage 26 Ingestion Results}
Evaluating the repeatable ingestion pipeline (`\texttt{STAGE26\_REPEATABLE\_DATA\_PIPELINE\_REPORT.md}`) resulted in an **8 / 8 PASS** result in Pytest (`\texttt{test\_pipeline\_stage26.py}`).

Key verification evidence:
\begin{itemize}
    \item \textbf{Schema Contract Enforcement}: Validated numeric rating bounds $[0..5]$ and trust score bounds $[0..100]$.
    \item \textbf{Dry-Run Safety}: `\texttt{pipeline.py dry-run}` generated `\texttt{dry\_run.json}` diff reports with **ZERO** PostgreSQL mutation.
    \item \textbf{Approval Protection}: Apply execution strictly rejected stale or invalid `\texttt{RUN\_ID}` parameters.
    \item \textbf{Selective Vector Synchronization}: Recalculate embeddings *only* for modified content hashes, skipping redundant vector operations.
\end{itemize}

\section{Stage 28 Master Orchestration Results}
Evaluating the single-command CLI orchestrator (`\texttt{STAGE28\_FINAL_REPORT.md}`) passed **6 / 6 tests** (`\texttt{test\_stage28\_orchestrator.py}`).

Executing `\texttt{python -m scripts.orchestrator full}` sequentially triggered all 6 upstream stages, generated `\texttt{final\_hotel\_dataset.csv}` (1,661 rows, 26 features), and created `\texttt{dry\_run.json}` under `\texttt{data/runs/<RUN\_ID>/}` while maintaining PostgreSQL read-only safety. Executing `\texttt{apply --run-id <RUN\_ID>}` verified dataset SHA-256 hashes before opening SQL transactions.

\section{Stage 29 Live Progress Results}
Stage 29 verification (`\texttt{STAGE29\_LIVE\_PROGRESS\_REPORT.md}`) passed **4 / 4 tests** (`\texttt{test\_stage29\_progress.py}`).

The `\texttt{ProgressTracker}` engine successfully rendered an interactive ASCII terminal progress dashboard showing active stage names, record percentage bars, elapsed execution time, and ETA calculations, outputting structured logs to `\texttt{pipeline.log}`. Emitting a `\texttt{SIGINT}` (Ctrl+C) signal cleanly terminated active sub-processes, marked `\texttt{pipeline\_manifest.json}` as `\texttt{INTERRUPTED}`, and guaranteed zero database corruption.

\section{Master Backend Automated Testing Results}
The entire backend test suite achieved a **50 / 50 PASSED (100\%)** result across all Pytest suites under tested conditions:
\begin{itemize}
    \item **Stage 29 Progress Suite**: **4 / 4 PASSED** (`\texttt{test\_stage29\_progress.py}`)
    \item **Stage 28 Orchestration Suite**: **6 / 6 PASSED** (`\texttt{test\_stage28\_orchestrator.py}`)
    \item **Stage 26 Ingestion Suite**: **8 / 8 PASSED** (`\texttt{test\_pipeline\_stage26.py}`)
    \item **Stage 24.5 Provenance Suite**: **18 / 18 PASSED** (`\texttt{test\_stage24\_5\_complete\_backend.py}`)
    \item **API, Context \& Grounding Suite**: **14 / 14 PASSED** (`\texttt{test\_api\_endpoints.py}`, etc.)
\end{itemize}

\section{Jupyter Research Notebook Experimental Findings}
The 10 Jupyter research notebooks (`\texttt{research/notebooks/}`) provided essential analytical evidence throughout system development:
\begin{enumerate}
    \item `\texttt{01\_hotel\_metadata\_analysis.ipynb}`: Mapped geospatial hotel clustering (New Delhi 216, Mahipalpur 56, Gurugram 39) and proved `\texttt{price\_level}` was 100\% missing (NaN).
    \item `\texttt{02\_review\_analysis.ipynb}`: Identified review positivity skew (median rating 4.10) and quantified the 5-review API detail cap limit.
    \item `\texttt{03\_sentiment\_analysis.ipynb}`: Verified DistilBERT sentiment polarity correlation ($r \approx 0.84$) against star ratings.
    \item `\texttt{04\_absa\_explainability\_analysis.ipynb}`: Evaluated 5 ABSA aspect distributions, identifying Cleanliness as the highest variance aspect differentiator.
    \item `\texttt{05\_feature\_engineering\_analysis.ipynb}`: Confirmed Trust Score Gaussian distribution (mean 68.0), Popularity power-law distribution, and Trust vs Popularity independence ($r \approx 0.05$).
    \item `\texttt{06\_user\_dataset\_analysis.ipynb}`: Audited demographic preference distributions across 500 synthetic user profiles.
    \item `\texttt{07\_interaction\_analysis.ipynb}`: Exposed 99.27\% matrix sparsity in initial V1 interactions, triggering interaction generator redesign.
    \item `\texttt{08\_final\_dataset\_overview.ipynb}`: Verified master dataset integrity (1,661 hotels, 26 features, zero nulls post-imputation).
    \item `\texttt{09\_recommender\_diagnostics.ipynb}`: Diagnosed Stage A model failure ($\text{NDCG}@10 = 0.006$, SVD underfitting, linear blending collapse).
    \item `\texttt{10\_interaction\_quality\_audit.ipynb}`: Audited V2 interactions, confirming preference overlap injection (66\% budget, 51\% area) and RRF performance recovery.
\end{enumerate}

\section{End-to-End System Results}
In the final production backend, executing user requests through FastAPI REST endpoints (`/api/v1/recommend`, `/api/v1/chat`, `/api/v1/hotel/{id}/explanation`) resolved queries against PostgreSQL 17.6 + `\texttt{pgvector}` in **210.0 ms cache miss latency** and **12.0 ms cache hit latency**.

\section{Limitations and Experimental Caveats}
The system's empirical findings operate under the following documented constraints:
\begin{itemize}
    \item **Google Places API Review Cap**: External API detail requests cap reviews at 5 per property, limiting long-tail review volume analysis.
    \item \textbf{Missing Price Attribute}: Google Places API returned 100\% missing `\texttt{price\_level}` data, requiring engineered \texttt{budget\_category} proxies.
    \item \textbf{Synthetic Interaction Modeling}: User interaction matrices were synthetically generated (V2 preference overlap model) rather than captured from live commercial transaction logs.
    \item \textbf{Local Execution Scope}: Local LLM inference (\texttt{mistral}/\texttt{llama3}) and PostgreSQL database deployments were evaluated on single-host developer infrastructure.
\end{itemize}

\section{Final Evaluation Summary}
Table~\ref{tab:final_evaluation_summary} summarizes the verified evaluation results across all system components.

\begin{table}[!ht]
\centering
\caption{Final Verification Summary across TrustLayer-AI System Components}
\label{tab:final_evaluation_summary}
\begin{tabular}{llll}
\hline
\textbf{Component} & \textbf{Evaluation Method} & \textbf{Verified Empirical Result} & \textbf{Status} \\
\hline
Data Pipeline & Dataset Integrity Audit & 1,661 Hotels, 7,910 Chunks, SHA-256 Verified & PASS \\
Sentiment / ABSA & DistilBERT & $r \approx 0.84$ Rating Correlation; Cleanliness Differentiator & PASS \\
Initial Recommender & Offline Benchmark (Stage A) & $\text{NDCG}@10 = 0.006$; SVD Underfitting ($\alpha=1.0$) & FAIL \\
Remediated Recommender & RRF Rank Fusion (Stage A.1) & $\text{NDCG}@10 > 0.12$; Calibration Mismatch Resolved & PASS \\
Explainability & Analytical Explainer & $< 5$ ms Latency; 100\% Readability Audit Pass & PASS \\
RAG Retrieval & 150-Query Benchmark & Precision@5 = 0.81, Recall@5 = 0.79, MRR = 0.83 & PASS \\
Grounding Validator & 150-Query Audit & 96.7\% Grounded Rate, 1.3\% Hallucination Rate & PASS \\
PostgreSQL / pgvector & Dual-Backend Parity & 1.0000 Cosine Sim Parity, 100\% Top-1 RAG Match & PASS \\
Repeatable Ingestion & Stage 26 Pytest Suite & 8 / 8 PASSED; Dry-Run Safety Verified & PASS \\
Master Orchestrator & Stage 28 Pytest Suite & 6 / 6 PASSED; One-Command CLI Verified & PASS \\
Progress & Stage 29 Pytest Suite & 4 / 4 PASSED; Live ASCII \& Ctrl+C Safety Verified & PASS \\
Master Backend Suite & Complete Pytest Suite & 50 / 50 PASSED (100\% Suite Verification) & PASS \\
\hline
\end{tabular}
\end{table}
